\section{Literature and Context}


\begin{frame}[allowframebreaks]{Literature and Context}
  \begin{block}{抵押品的作用}
    \begin{itemize}
      \item \textcite{stiglitz1981credit} 提出了\textbf{信贷配给问题}：在信息不对称下，提高利率会恶化借款人类型与行为，银行因此可能宁可拒绝部分愿意按当前利率借款的申请者，也不愿继续提价。
      \item \textcite{bester1985screening} 的研究表明，这一问题可以通过将抵押品作筛选工具来(部分)解决：信用风险较低的借款人提供更多抵押品并获得较低利率，从而消除了配给的必要性。
      \item 使用抵押品的另一种解释是为了缓解道德风险问题：提供抵押品增加了借款人进行风险转移、逃避责任或策略性违约的成本 \parencite{bester1987role,chan1987collateral,tirole2006theory}。
    \end{itemize}
  \end{block}


  \medskip
  \begin{block}{发展中国家的教育}

    \begin{itemize}
      \item \textcite{williams2015education} 指出，在不少非洲国家，家庭难以承担学费、书本费、校服费、交通费和餐饮费等教育自付支出。
      \item \textcite{duflo2021impact} 表明，为经济困难学生提供奖学金，不仅提高了中等和高等教育完成率，还改善了长期劳动力市场表现。
      \item 贷款在高等教育中更为常见，并已在智利、南非和中国等中等收入国家得到研究 \parencite{bucarey2020labor,gurgand2023student,cai2019repayment}。与此同时，也有研究强调传统学生贷款安排可能带来较重的偿还负担与违约风险 \parencite{cai2019repayment,dearden2019evaluating}。
    \end{itemize}
    
  \end{block}
\end{frame}
